\documentclass[conference,a4paper]{IEEEtran}
\IEEEoverridecommandlockouts
\usepackage{cite}
\usepackage{amsmath,amssymb,amsfonts}
\usepackage{graphicx}
\usepackage{textcomp}
\usepackage{xcolor}
\usepackage{booktabs}
\usepackage{url}

\begin{document}

\title{What Moves a Quantum Advantage? Initialisation, Training Protocol and
Circuit Ablation in Hybrid Quantum-Classical PINNs}

\author{
\IEEEauthorblockN{J Abishek Rufus Raj}
\IEEEauthorblockA{\textit{CHRIST (Deemed to be University), Kengeri Campus} \\
Bengaluru, India \\
abishekrufusraj@gmail.com}
\and
\IEEEauthorblockN{Mary Anita E A}
\IEEEauthorblockA{\textit{CHRIST University} \\
Bengaluru, India \\
maryanita.ea@christuniversity.in}
}

\maketitle

\begin{abstract}
Hybrid quantum-classical physics-informed neural networks (PINNs) are reported to match
classical accuracy at a fraction of the trainable parameters. We ask what such a
comparison actually measures. Using a published quantum PINN's own released
implementation, unmodified, on its own two-dimensional Helmholtz problem, we reproduce
its result: a median relative $L_2$ error of 5.16\% over ten seeds against a published
$6.69 \pm 1.51\%$, beating its classical baseline at 9.88\% (exact permutation test,
$p = 0.0044$, Cliff's $\delta = -0.84$). We then vary two implementation choices that the
comparison holds fixed. Initialising the classical baseline as the paper's text describes,
rather than as its code does, halves the effect size to $\delta = -0.50$. Training the
same baseline to a convergence criterion instead of a fixed epoch budget takes it to
0.70\%, a factor of seventeen, where the advantage claimed is a factor of about two.
Under the original protocol, a Xavier-initialised classical PINN already matches the
quantum model (4.63\% against 5.16\%, $p = 0.32$). A fully symmetric control --- both arms
from the authors' code, one protocol, ten seeds --- puts the quantum model ahead of
\emph{their} classical solver (1.73\% against 4.30\%, $p = 0.0068$), yet our own
implementation of that same 2{,}751-parameter network reaches 0.70\% under the identical
criterion, so the quantum lead measures the baseline's implementation and not its
architecture. Ablation locates what the circuit
contributes: in our own architecture, replacing it with a classical map of the same shape
changes nothing ($p = 0.79$ and $p = 0.62$ for two ansaetze, eight seeds each), while in
the published architecture the circuit is load-bearing but works frozen at initialisation
(7.84\%, no failures). We conclude that reported quantum advantages in this setting are
sensitive to baseline initialisation and training protocol, and that neither can be left
unreported.
\end{abstract}

\begin{IEEEkeywords}
quantum machine learning, physics-informed neural networks, variational quantum circuits,
reproducibility, ablation, experimental protocol
\end{IEEEkeywords}

\section{Introduction}

A physics-informed neural network (PINN) trains a network to satisfy a differential
equation at sampled collocation points, using the equation itself as the training
signal~\cite{raissi}. A hybrid quantum-classical PINN replaces part of that network with a
variational quantum circuit. The motivation is representational: an angle-encoded circuit
computes a truncated Fourier series whose accessible frequencies are fixed by its encoding
gates rather than by its trainable weights~\cite{schuld}, so a circuit might express
oscillatory solution structure with fewer parameters than a classical layer.

Reported results support this. A recent hybrid quantum-classical PINN obtains comparable
accuracy using 10--30\% of the trainable parameters of a classical
baseline~\cite{qcpinn}. Such comparisons, however, fix a number of implementation choices
that are rarely varied and sometimes not reported: how each arm is initialised, how long
each arm is trained, and how the stopping point is decided. If those choices move the
result by as much as the architecture does, the comparison does not measure what it is
taken to measure.

This paper asks how much they move it. Our contributions are:

\begin{enumerate}
\item A replication of a published quantum PINN comparison using the authors' released
implementation, unmodified, on their own problem, with recorded seeds. The published
result reproduces (Section~\ref{sec:rep}).
\item A measurement of how far that comparison moves under two choices it holds fixed:
the baseline's initialisation, which halves the effect size, and the training protocol,
which moves the baseline by a factor of seventeen
(Sections~\ref{sec:init}--\ref{sec:proto}).
\item A symmetric control that trains both arms under one protocol, and the implementation
gap it exposes in the published classical baseline (Section~\ref{sec:sym}).
\item An ablation, with a frozen-circuit control, separating what the circuit contributes
in two different hybrid architectures (Section~\ref{sec:abl}).
\end{enumerate}

All comparisons use repeated seeds and exact permutation tests; effect sizes are reported
alongside $p$-values throughout.

\section{Related Work}

The expressivity argument for quantum layers rests on the encoding-determined frequency
spectrum of variational circuits~\cite{schuld}. Against it stands trainability: gradient
variance in random parameterised circuits vanishes exponentially in qubit
count~\cite{mcclean}, an effect that depends on whether the cost function is global or
local~\cite{cerezo}, and which is commonly mitigated by initialising the circuit close to
the identity~\cite{grant}. Expressibility and entangling capability~\cite{sim} and
dimensional expressivity analysis~\cite{funcke} provide standard vocabulary for comparing
ansaetze, and we use both.

On the classical side, PINN training failures are optimisation failures rather than
capacity failures~\cite{krishnapriyan}, and a line of work on learned or enhanced
initialisation shows that initial weights drive catastrophic convergence failures in
classical PINNs. Benchmarking studies of quantum machine learning outside the PINN setting
report that classical models frequently match or exceed quantum ones, and that removing
entanglement often leaves performance unchanged~\cite{bowles}. Our ablation applies that
style of control to a published quantum PINN rather than to classifiers.

\section{Experimental Protocol}

\textbf{Problems.} Two settings are used. For the replication we use the
two-dimensional Helmholtz problem of~\cite{qcpinn}: $u_{x_1x_1} + u_{x_2x_2} + \lambda u
= f$ on $[-1,1]^2$ with $\lambda = 1$, exact solution $u^\star = \sin(\pi
x_1)\sin(4\pi x_2)$ and homogeneous Dirichlet boundaries. For the ablation of our own
architecture we use the viscous Burgers' equation with $\nu = 0.01/\pi$, the standard PINN
stress test.

\textbf{Architectures.} The replicated model places a five-qubit angle-cascade circuit
(15 quantum parameters) between two small classical blocks, 771 parameters in total, and
is compared against a 2{,}751-parameter classical network~\cite{qcpinn}. Our own
architecture instead places a four-qubit circuit at the input, feeding a seven-layer
classical stack; two entangling blocks are used, a ring of $R_X$ rotations with CNOTs (8
quantum parameters) and a strongly entangling block of full rotations (24), giving 2{,}649
and 2{,}665 parameters against a 3{,}021-parameter classical PINN.

\textbf{Training.} Two protocols are compared explicitly. \emph{Theirs} is the published
recipe: Adam at $5\times10^{-3}$, mini-batches of 64, a fixed budget of 20{,}000 epochs
with a plateau scheduler, and no second-order phase. \emph{Ours} terminates on a
convergence criterion --- an exponential moving average of the loss improving by less than
3\% over 3{,}000 epochs --- followed by 500 L-BFGS steps, a two-stage schedule standard in
the PINN literature~\cite{raissi}.

\textbf{Statistics.} Sample sizes are small, so all tests are exact: two-sided permutation
tests on the difference in medians, enumerated completely, with the observed labelling
counted in the numerator. Effect sizes are Cliff's $\delta$~\cite{cliff}. Ten seeds per arm are used in
the replication, matching the source; eight seeds per arm elsewhere. Seeds are recorded;
the released implementation we replicate carries a seed field that its code never applies,
so its published runs cannot be reproduced individually.

\begin{table*}[t]
\caption{Full configuration of both hybrid architectures. The replicated column
describes the released implementation of~\cite{qcpinn} as run here; the left column
describes this work. Parameter counts are read from the trained models.}
\label{tab:config}
\centering
\scriptsize
\begin{tabular}{p{0.17\textwidth}p{0.38\textwidth}p{0.36\textwidth}}
\hline
& \textbf{This work (circuit at the input)} & \textbf{Replicated~\cite{qcpinn} (circuit in the middle)} \\
\hline
Qubits & 4 & 5 \\
Circuit depth & 2 data-reuploading layers, entangler depth 1 & 1 quantum layer \\
Encoding & angle, $R_X$; the pair $(x,t)$ tiled across the 4 wires, re-encoded before each layer & angle, $R_X$, applied once \\
Ansatz & \emph{ring}: one $R_X$ per wire then a CNOT ring (\textsc{BasicEntanglerLayers}); \emph{strong}: $R_ZR_YR_Z$ per wire then CNOTs (\textsc{StronglyEntanglingLayers}) & \emph{cascade}: one $R_X$ and one $R_Z$ per wire, then a ring of CRX gates \\
Observables & $\langle Z_i\rangle$ on each of the 4 wires & $\langle Z_i\rangle$ on each of the 5 wires \\
Classical blocks & none before; 7 hidden layers of width 20 after & $2\!\to\!50\!\to\!5$ before, $5\!\to\!50\!\to\!1$ after \\
Quantum parameters & 8 (ring), 24 (strong) & 15 \\
Total parameters & 2{,}649 (ring), 2{,}665 (strong); classical control 3{,}021 & 771; classical baseline 2{,}751 \\
Quantum initialisation & $\mathcal{U}[0,2\pi)$ (library default) or $\mathcal{N}(0,0.1)$ (near-identity~\cite{grant}) & Xavier~\cite{glorot}, as stated by the source \\
Classical initialisation & Xavier~\cite{glorot} & Xavier~\cite{glorot} \\
Optimiser & Adam, $10^{-3}$, then 500 L-BFGS steps (strong-Wolfe) & Adam, $5\times10^{-3}$, \textsc{ReduceLROnPlateau}, gradient clipping at 1.0 \\
Training points & 1{,}400 per epoch (800 interior, 300 shock-refined, 150 initial, 150 boundary), redrawn every epoch & mini-batches of 64, redrawn every epoch \\
Loss weights & residual $+\,10\times$ (initial $+$ boundary) & residual $+\,10\times$ boundary \\
Epochs & convergence criterion (EMA improving by $<3\%$ over 3{,}000 epochs), cap 25{,}000 & fixed 20{,}000 \\
Simulation & PennyLane 0.44 \texttt{default.qubit}, analytic (no shots), backpropagation & PennyLane 0.29 \texttt{default.qubit}, analytic (no shots) \\
\hline
\end{tabular}
\end{table*}

\textbf{Simulation environment.} All runs are noiseless statevector simulations on CPU.
This work uses PennyLane~\cite{pennylane} 0.44 with PyTorch 2.10 on Python 3.13; the replication runs the
authors' released code unmodified in its own environment, PennyLane 0.29 with PyTorch 2.0.1
on Python 3.9, at commit \texttt{0dc467be}. Both execute on one 8-core Apple M2 with 16\,GB
of memory, two threads per run, four runs concurrently. No quantum hardware is used, and
no shot noise is simulated: every expectation value is exact.

\section{Replicating a Published Comparison}
\label{sec:rep}

We ran the authors' public implementation~\cite{qcpinn} (MIT licence) without
modification, in their pinned package versions, on their Helmholtz problem, adding only a
per-run seed. Fig.~\ref{fig:one} shows every run.

\begin{figure*}[t]
\centering
\includegraphics[width=\linewidth]{fig1.pdf}
\caption{Per-seed relative $L_2$ error on the Helmholtz problem; every point is one
training run, bars are medians, and the axis is logarithmic. The published comparison is
the first two columns. Columns three to five vary one choice at a time: the baseline's
initialisation, the baseline trained to a convergence criterion, and a classical PINN
under the published protocol. The two columns right of the rule are the symmetric control
of Section~\ref{sec:sym}, where both arms come from the authors' code and differ by
nothing but the model; the classical arm's three runs near 100\% are dead
initialisations.}
\label{fig:one}
\end{figure*}

The published result reproduces. The quantum model reaches a median 5.16\% against a
published $6.69 \pm 1.51\%$, solves the problem on all ten seeds, and beats its classical
baseline at 9.88\% ($p = 0.0044$, $\delta = -0.84$). We find no discrepancy with the
reported measurement.

Two properties of the comparison are worth recording. First, the baseline is unstable:
three of its ten seeds fail outright, and those failures carry the mean. Second, the
efficiency is in parameters, not compute. A quantum run takes 78 minutes on a CPU against
one minute for the classical baseline, so a $3.5\times$ parameter saving costs roughly
$80\times$ the energy in simulation.

\section{Initialisation Inside the Baseline}
\label{sec:init}

The paper states that its classical layers use Xavier initialisation. Its released code
applies Xavier to the pre-processor only, leaving the remaining layers at framework
defaults. We therefore re-ran the baseline with Xavier applied to every classical layer,
as the text describes, changing nothing else.

\begin{table}[t]
\caption{The published comparison under two initialisations of its classical baseline.
Ten seeds per arm.}
\label{tab:init}
\centering
\footnotesize
\begin{tabular}{lccc}
\toprule
Baseline init. & Median & Fails & $\delta$ vs.\ quantum \\
\midrule
as released (first layer) & 9.88\% & 3/10 & $-0.84$ ($p{=}0.0044$) \\
as described (all layers) & 8.47\% & 1/10 & $-0.50$ ($p{=}0.0305$) \\
\bottomrule
\end{tabular}
\end{table}

Table~\ref{tab:init} shows the effect. One initialisation line, internal to one arm,
halves the effect size and removes two of the three failures. The quantum model remains
ahead, but the magnitude of the advantage is a function of a choice the comparison does
not vary.

This is consistent with what initialisation does to a purely classical PINN under our own
protocol on Burgers': rescaling the initial weights to a Xavier-equivalent gain changes the
seed-to-seed standard deviation by a factor of 23.1 and takes the failure count from three
in eight to zero. Initialisation is not a detail of implementation in this setting; it is
an experimental factor of the same order as the architectural differences being reported.

\section{Training Protocol}
\label{sec:proto}

The published protocol stops at a fixed budget. We retrained the same classical baseline
--- same architecture, same parameter count, same problem --- under a convergence criterion
with an L-BFGS phase. It reaches a median 0.70\% (Fig.~\ref{fig:one}), against the 12.12\%
reported and the 9.88\% we reproduce.

The baseline was therefore evaluated roughly a factor of seventeen from its own
capability, while the advantage being claimed is a factor of about two. We also ran the
comparison the other way: under the \emph{published} protocol, a Xavier-initialised
classical PINN reaches 4.63\% against the quantum model's 5.16\%, a difference that is not
established at eight seeds ($p = 0.3158$). The quantum advantage does not survive fixing
either the baseline's initialisation or its training.

What follows from these measurements is that the published baseline was trained far short
of its capability. Whether the advantage survives training both arms alike is a separate
question, and Section~\ref{sec:sym} answers it directly.

\section{A Symmetric Control}
\label{sec:sym}

Sections~\ref{sec:init}--\ref{sec:proto} vary one arm at a time, so the standing objection
is that no single experiment trains both arms alike. We ran that control directly: the
authors' quantum solver and the authors' classical solver, from the same unmodified
release, under one protocol --- the convergence criterion of Section~\ref{sec:proto}, a
500-step L-BFGS phase refined on 1{,}400 points, identical evaluation grid, seeds 0--9.
Nothing differs between the arms except the model.

The quantum arm reaches a median 1.73\% relative $L_2$ error (IQR 1.61--2.01, range
1.21--2.31) with no failures in ten seeds. The classical arm reaches a median 4.30\%, but
its distribution is bimodal: seven seeds land between 1.21\% and 7.02\%, and three (seeds
5, 6 and 9) finish near 100\% with a residual error of 99.99998\%, meaning the network
output is identically zero. Those same three seeds fail identically under the published
protocol, so they are dead initialisations rather than an artefact of the stopping rule,
and we report rather than discard them.

Over all ten seeds the difference is significant: median $-2.57$ pp, exact permutation
$p = 0.0068$ over all $\binom{20}{10} = 184{,}756$ splits, Cliff's $\delta = -0.64$. That
result is carried substantially by the three dead seeds. Among the seven seeds where the
classical arm trains at all the medians are 1.77\% against 3.43\%, $\delta = -0.43$, and
the difference is not established ($p = 0.0583$); the failure counts alone, 0/10 against
3/10, are not established either (Fisher $p = 0.2105$). The defensible statement is that
under a symmetric protocol the quantum model is more \emph{reliable} than this classical
arm, and that its accuracy margin over the seeds that train is a medium effect ten seeds
cannot resolve.

It does not follow that the circuit beats the architecture. The same 2{,}751-parameter
classical network, on the same equation under the same convergence criterion and L-BFGS
phase, but implemented in our own trainer with a 1{,}400-point collocation budget, reaches
a median 0.70\% over eight seeds (range 0.53--0.84) with no failures: better than the
quantum arm by a factor of 2.5, and better than the authors' own classical solver by a
factor of six. The quantum lead in the symmetric control is therefore a property of the
classical baseline \emph{as implemented}, not of the classical architecture --- the same
conclusion Sections~\ref{sec:init}--\ref{sec:proto} reach from the other direction. The
quantum arm also costs a median 51.7 minutes against 0.86, a factor of 60
(Section~\ref{sec:cost}).

\textbf{Does the lead survive other training settings?} The symmetric protocol fixes one
optimiser, learning rate and batch size, and a lead that held only at that point in
hyperparameter space would be an artefact of it. We therefore re-ran both arms under five
further settings --- SGD with momentum, RMSprop, learning rates of $10^{-3}$ and
$10^{-2}$, and a batch of 256 --- at three seeds each, thirty runs in all. The quantum arm
reaches the lower median in every setting where median and mean agree on the ordering,
five of six; RMSprop splits, because one classical seed there fails outright at 100.7\%
and so dominates that cell's mean. Three seeds cannot establish any single comparison ---
the exact two-sided permutation floor is $2/\binom{6}{3} = 0.10$ --- so we report
direction rather than per-cell significance, and the sign test over the five agreeing
cells returns $p = 0.0625$, itself the smallest value five cells admit. Pooled across all
six settings the failure asymmetry of this section reappears where we were not looking for
it: the quantum arm's worst run of eighteen is 9.71\%, the classical arm's worst is
100.7\%, and dead initialisations run 0/18 against 2/18 (Fisher $p = 0.4857$). The
direction is consistent across every setting we varied; its magnitude is not resolved at
this seed count.

\section{What Does the Circuit Contribute?}
\label{sec:abl}

The direct test of whether a quantum layer earns its place is to remove it and put a
classical map of the same shape in the same slot~\cite{bowles}. Both stand-ins we use
carry slightly more parameters than the circuit they replace, so the classical side is
favoured.

\begin{figure}[t]
\centering
\includegraphics[width=\linewidth]{fig2.pdf}
\caption{Circuit ablation in two hybrid architectures. (a) With the circuit at the input
of a deep classical stack, removing it changes nothing. (b) With the circuit inside a
five-unit bottleneck, removing it is fatal, but freezing its weights at initialisation is
not. Note the logarithmic scale.}
\label{fig:two}
\end{figure}

The two architectures answer differently (Fig.~\ref{fig:two}). In our architecture, where
the circuit feeds a deep classical stack, replacing it with a linear map changes nothing
measurable: 3.73\% to 3.23\% for the ring ansatz ($p = 0.79$) and 3.33\% to 2.69\% for the
strongly entangling one ($p = 0.62$), with the ablated models marginally ahead. Every
quantum-assisted result we obtained in that architecture came from a model whose circuit
could have been deleted.

In the published architecture the circuit is load-bearing: both classical stand-ins fail
to train at all, never bringing the residual below $\sim\!5\times10^3$ where the circuit
brings it to about 2. A frozen-circuit control separates two explanations. Keeping the
circuit but freezing its fifteen weights at initialisation still solves the problem on all
ten seeds at 7.84\%, against 5.16\% trained ($p = 0.0032$). The circuit's contribution
there is therefore mostly as a \emph{fixed} nonlinear map conditioning a narrow
bottleneck; training its quantum parameters is worth about 2.7 percentage points on top.

Two supporting measurements agree. Across circuit sizes from two to eight qubits, the
circuit receives between 0.06\% and 0.4\% of the gradient norm, the remainder going to the
classical weights. And the strongly entangling ansatz carries redundant parameters by
construction: probing the Jacobian of its output gives about nine effective directions from
24 parameters, with its final $R_Z$ unobservable before a Pauli-$Z$ measurement, the
kind of superfluous parameter dimensional expressivity analysis is designed to
detect~\cite{funcke}.

\section{Architecture and Initialisation Interact}

If initialisation merely shifted every architecture equally, comparisons at one fixed
choice would still rank architectures correctly. It does not. On the Helmholtz problem,
with eight seeds per arm under our protocol, near-identity initialisation leaves the
strongly entangling ansatz at 7.54\% and takes the ring ansatz to 66.06\%, failing on
seven seeds of eight. The ansatz contrast under that initialisation is the largest effect
we measure ($\delta = 0.94$, $p = 0.0031$, Holm-adjusted~\cite{holm} 0.019 within a family of six).
The same interaction, in the same direction, appears on Burgers'.

The mechanism follows from the gate algebra. The ring's trainable gates are $R_X$
rotations, the same axis as the encoding, so $R_X(w)R_X(f) = R_X(f+w)$ and its weights act
only as additive offsets to the encoding angles; it can generate no rotation about $Y$ or
$Z$. Near-identity initialisation sets those offsets near zero. The strongly entangling
block, rotating about all three axes, is not so constrained.

\section{A Third Equation, and the Circuit's Components}
\label{sec:third}

Two equations cannot separate a property of hybrid PINNs from a property of the
Helmholtz problem, so we add a third: the one-dimensional heat equation, under the
convergence protocol of Section~\ref{sec:proto}, this work's architecture, eight seeds per
arm. The ordering of the other two equations reproduces. The classical network reaches a
median 0.33\%, the strong ansatz 0.48\% and the ring ansatz 0.67\%, in that order on the
mean as well as the median. The classical arm beats the ring circuit at the smallest
$p$ this design admits ($p = 40/12{,}870 = 0.0031$, $\delta = +0.88$) and leads the strong
circuit without reaching significance ($p = 0.0796$, $\delta = +0.66$). Three equations,
one ranking.

\textbf{Depth.} Doubling the variational depth from two layers to four buys nothing on
Burgers'. The strong ansatz is unmoved --- 3.33\% to 3.43\% ($p = 0.9378$) --- for twice
the circuit parameters, and the ring ansatz does not improve either, its median and mean
both moving the wrong way (3.73\% to 9.44\%, $p = 0.1064$). Whatever the circuit
contributes here, more of it does not contribute more.

\textbf{Entanglement structure.} The strong ansatz carries three rotation axes per wire
and twenty-four circuit parameters against the ring's single axis and eight. It reaches
the lower median on all three equations --- 3.33\% against 3.73\% on Burgers', 8.90\%
against 17.99\% on Helmholtz, 0.48\% against 0.67\% on heat --- but no single comparison
is established at eight seeds ($p = 0.79$, $0.067$ and $0.13$). The direction is
consistent across three equations; the magnitude is not resolved by this many seeds.

\textbf{Qubit count.} We do not vary it at a seed count that would support a claim. The
runs we hold at other circuit widths are single-seed, and this study's standard is eight;
we therefore report no number for them. This is the one circuit component we leave open,
and the released configurations are sufficient to run it.

\section{Costs: Parameters, Wall-Clock and Energy}
\label{sec:cost}

Reports in this area often state a parameter saving and leave the reader to assume it is a
saving. The three costs move in opposite directions, so we separate them.

\textbf{Parameters.} The saving is real. The replicated hybrid reaches 5.16\% with 771
parameters where its classical baseline needs 2{,}751 for 9.88\%, a $3.6\times$ reduction
at better accuracy under the published protocol. In our own architecture the saving is
negligible: the circuit replaces so little of the network that 2{,}649 parameters stand
against a 3{,}021-parameter control, $1.14\times$.

\textbf{Wall-clock.} The saving inverts, by two orders of magnitude.
Table~\ref{tab:cost} gives medians over the seeds of each arm. Simulating a 5-qubit
circuit for 20{,}000 epochs costs 71.0 minutes against 0.8 minutes for the classical
baseline, $89\times$; our 4-qubit circuit costs $17\times$ its classical control on
Burgers' and $15\times$ on Helmholtz. The circuit ablations locate that cost precisely:
replacing the circuit with a classical map of the same width leaves accuracy unchanged on
Burgers' (3.73\% to 3.23\%) while returning the run to 2.6 minutes from 34.3.

\textbf{Energy.} We did not meter power, so we report no absolute figures. Every arm ran
on the same processor at the same thread count and saturated it, so to the accuracy of
that assumption energy is proportional to wall-clock and the ratios above carry over. We
state this as an inference, not a measurement.

\textbf{What a hardware cost would be instead.} None of these wall-clock figures is a
quantum runtime, and they should not be read as predictions of one. Simulation cost grows
exponentially in qubit count, which is the reason the figures above are large; hardware
cost does not, and would instead be dominated by circuit executions, shot noise and queue
latency, none of which appear here. A 4- or 5-qubit circuit of depth 1--2 is small enough
that per-shot execution would be fast, but obtaining the gradients these models train on
would require many thousands of circuit executions per epoch, and the analytic
expectation values used throughout this study would have to be replaced by shot estimates.
Whether the parameter saving survives that substitution is an open question which our
measurements do not address.

\begin{table}[t]
\caption{Median cost and accuracy per arm, over 10 seeds for the replication and 8
elsewhere. Wall-clock is total time for the arm's own protocol, read from the run
artefacts. Quantum parameters are a subset of the total.}
\label{tab:cost}
\centering
\scriptsize
\setlength{\tabcolsep}{1.5pt}
\begin{tabular}{llrrr}
\hline
\textbf{Problem / arm} & \textbf{Circuit} & \textbf{Par.} & \textbf{Min.} & \textbf{Err.\%} \\
\hline
\multicolumn{5}{l}{\emph{Helmholtz, published protocol, 20{,}000 epochs}~\cite{qcpinn}} \\
\quad hybrid & 5 qubits, 15 q.par. & 771 & 71.0 & 5.16 \\
\quad classical baseline & --- & 2{,}751 & 0.8 & 9.88 \\
\quad hybrid, circuit frozen & 5 qubits, 0 trained & 771 & 58.7 & 7.84 \\
\quad hybrid, circuit removed & classical $\tanh$ map & 786 & 1.1 & 1342 \\
\hline
\multicolumn{5}{l}{\emph{Burgers', convergence protocol (this work)}} \\
\quad QAPINN ring & 4 qubits, 8 q.par. & 2{,}649 & 34.3 & 3.73 \\
\quad QAPINN strong & 4 qubits, 24 q.par. & 2{,}665 & 35.0 & 3.33 \\
\quad ring, circuit removed & classical map & 2{,}653 & 2.6 & 3.23 \\
\quad strong, circuit removed & classical map & 2{,}673 & 4.1 & 2.69 \\
\quad classical control & --- & 3{,}021 & 2.0 & 2.21 \\
\hline
\multicolumn{5}{l}{\emph{Helmholtz, convergence protocol (this work)}} \\
\quad QAPINN strong, near-identity & 4 qubits, 24 q.par. & 2{,}665 & 66.9 & 7.54 \\
\quad QAPINN ring, near-identity & 4 qubits, 8 q.par. & 2{,}649 & 52.9 & 66.06 \\
\quad classical control & --- & 3{,}021 & 4.6 & 6.90 \\
\hline
\end{tabular}
\end{table}

\textbf{Simulation is not hardware.} Every wall-clock and energy figure in this section is
CPU statevector simulation, and the quantum arm's cost here is the cost of simulating a
circuit, not of running one. The two scale in opposite directions: simulation grows
exponentially in qubit count, where hardware execution would not, so the factor of 60 in
Table~\ref{tab:cost} is an artefact of the method of evaluation and must not be read as a
hardware projection. What the simulated figures do establish is the cost of the
\emph{research} --- the price of obtaining these results at this width --- and the
parameter counts, which are hardware-independent. A hardware comparison would have to
account for shot noise, queueing and transpilation, none of which appear here, and we make
no claim about it.

\section{Limitations}

Our study is a simulation study: four- and five-qubit circuits, three PDEs, statevector
simulation on a CPU, no quantum hardware --- as in every comparable study we examined. The
symmetric control of Section~\ref{sec:sym} is confined to one equation and to the authors'
own implementation of both arms; we have not run it on Burgers' nor against our own
quantum models, and at ten seeds per arm the accuracy margin among non-failing seeds
remains a medium effect we cannot resolve. Our own
classical arm is also not the strongest classical model available: on Helmholtz the
replicated 2{,}751-parameter network reaches 0.70\% under our protocol where our deeper
network reaches 4.31\%, which makes our quantum-versus-classical comparisons conservative
rather than favourable. Finally, at eight seeds an exact median-difference test has an
attainable $p$-floor of $40/12{,}870 = 0.0031$; ten seeds per arm are required for such a
test to clear a correction over twenty comparisons, and we report where that bound binds.

\section{Conclusion}

A published quantum-classical PINN advantage reproduces exactly under its authors' own
code, and then moves by half its effect size when the classical baseline is initialised as
that paper describes, and inverts when the baseline is implemented and trained to its own capability.
Ablation shows the circuit contributes nothing measurable in one hybrid architecture and
acts as a fixed feature map in the other. We therefore recommend that comparisons in this
area report and vary both the initialisation of every arm and the stopping rule, and that
effect sizes be reported with repeated seeds rather than single runs. The code, the configuration files, the per-run seeds, the raw
run artefacts and a verification suite that re-derives every number in this paper from
them are released at \url{https://github.com/ANONYMISED/qapinn-systematic-study}, so that
any figure here can be checked without rerunning the training.

\begin{thebibliography}{00}
\bibitem{raissi} M. Raissi, P. Perdikaris, and G. E. Karniadakis, ``Physics-informed neural networks: a deep learning framework for solving forward and inverse problems involving nonlinear partial differential equations,'' \emph{J. Comput. Phys.}, vol. 378, pp. 686--707, 2019.
\bibitem{schuld} M. Schuld, R. Sweke, and J. J. Meyer, ``Effect of data encoding on the expressive power of variational quantum machine-learning models,'' \emph{Phys. Rev. A}, vol. 103, p. 032430, 2021.
\bibitem{qcpinn} A. Farea, S. Khan, and M. S. Celebi, ``QCPINN: quantum-classical physics-informed neural networks for solving PDEs,'' \emph{Mach. Learn.: Sci. Technol.}, 2025; arXiv:2503.16678v6.
\bibitem{mcclean} J. R. McClean, S. Boixo, V. N. Smelyanskiy, R. Babbush, and H. Neven, ``Barren plateaus in quantum neural network training landscapes,'' \emph{Nat. Commun.}, vol. 9, p. 4812, 2018.
\bibitem{cerezo} M. Cerezo, A. Sone, T. Volkoff, L. Cincio, and P. J. Coles, ``Cost-function-dependent barren plateaus in shallow parametrized quantum circuits,'' \emph{Nat. Commun.}, vol. 12, p. 1791, 2021.
\bibitem{grant} E. Grant, L. Wossnig, M. Ostaszewski, and M. Benedetti, ``An initialization strategy for addressing barren plateaus in parametrized quantum circuits,'' \emph{Quantum}, vol. 3, p. 214, 2019.
\bibitem{sim} S. Sim, P. D. Johnson, and A. Aspuru-Guzik, ``Expressibility and entangling capability of parameterized quantum circuits for hybrid quantum-classical algorithms,'' \emph{Adv. Quantum Technol.}, vol. 2, no. 12, p. 1900070, 2019.
\bibitem{funcke} L. Funcke, T. Hartung, K. Jansen, S. K\"uhn, and P. Stornati, ``Dimensional expressivity analysis of parametric quantum circuits,'' \emph{Quantum}, vol. 5, p. 422, 2021.
\bibitem{krishnapriyan} A. S. Krishnapriyan, A. Gholami, S. Zhe, R. M. Kirby, and M. W. Mahoney, ``Characterizing possible failure modes in physics-informed neural networks,'' in \emph{Adv. Neural Inf. Process. Syst.}, vol. 34, 2021.
\bibitem{bowles} J. Bowles, S. Ahmed, and M. Schuld, ``Better than classical? The subtle art of benchmarking quantum machine learning models,'' arXiv:2403.07059, 2024.
\bibitem{glorot} X. Glorot and Y. Bengio, ``Understanding the difficulty of training deep feedforward neural networks,'' in \emph{Proc. AISTATS}, pp. 249--256, 2010.
\bibitem{pennylane} V. Bergholm et al., ``PennyLane: automatic differentiation of hybrid quantum-classical computations,'' arXiv:1811.04968, 2018.
\bibitem{holm} S. Holm, ``A simple sequentially rejective multiple test procedure,'' \emph{Scand. J. Stat.}, vol. 6, pp. 65--70, 1979.
\bibitem{cliff} N. Cliff, ``Dominance statistics: ordinal analyses to answer ordinal questions,'' \emph{Psychol. Bull.}, vol. 114, no. 3, pp. 494--509, 1993.
\end{thebibliography}

\end{document}
